\chapter{Introdução}
%\addcontentsline{toc}{chapter}{\uppercase{INTRODUÇÃO}}
\label{cap:introducao}
%\phantom{0}

O câncer de próstata é uma das neoplasias malignas mais comuns em homens, representando um desafio significativo para a saúde pública~\cite{ref-siegel2024}. A graduação tumoral desempenha um papel central na definição do manejo clínico e do prognóstico. O sistema de Gleason, baseado na soma dos dois padrões histológicos mais prevalentes (escore de Gleason), evoluiu para o sistema de graduação da ISUP (0–5), que fornece uma estratificação prognóstica aprimorada e é amplamente utilizado na prática clínica~\cite{ref-epstein2016}.

A avaliação histopatológica tradicional de imagens de lâminas inteiras (\textit{Whole-Slide Images} – WSIs) depende da análise manual realizada por patologistas, um processo sujeito à variabilidade inter e intraobservador~\cite{ref-bulten2020}, além de ser demorado e limitado pela disponibilidade de especialistas. Essas restrições motivaram o desenvolvimento de métodos automatizados baseados em aprendizado profundo.

Redes neurais convolucionais (CNNs) têm demonstrado desempenho destacado na análise de imagens médicas~\cite{ref-campanella2019}, incluindo a graduação automatizada de padrões de Gleason e graus ISUP~\cite{ref-nagpal2019,ref-bulten2020}. Entre as arquiteturas modernas, a família EfficientNet~\cite{ref-tan2019} oferece um excelente equilíbrio entre acurácia e custo computacional, devido à sua estratégia de \textit{compound scaling} e ao uso de blocos MBConv com \textit{squeeze-and-excitation}~\cite{ref-hu2018}.

Neste trabalho, investigamos a classificação automatizada dos graus ISUP por meio da avaliação de variantes da EfficientNet (B0, B3 e B7) no conjunto de dados PANDA~\cite{ref-panda2020}. Empregamos técnicas de \textit{transfer learning}, estratégias de \textit{fine-tuning}, técnicas avançadas de treinamento e métodos de \textit{ensemble} para aumentar a robustez dos modelos. 


\section{Objetivo Geral}
\label{sec:objetivo-geral}

Investigar e avaliar métodos de aprendizado profundo para a classificação automatizada dos graus ISUP em câncer de próstata, utilizando variantes da arquitetura EfficientNet, com foco no aprimoramento da robustez, da estabilidade e do desempenho dos modelos.

\section{Objetivos Específicos}
\label{sec:objetivos-especificos}

Para alcançar o objetivo geral deste trabalho, faz-se necessário atingir os seguintes objetivos específicos:

\begin{itemize}
    \item Implementar e avaliar uma função de perda que combine componentes ordinais e focais, visando melhorar o desempenho em tarefas de classificação ordinal.
    \item  Comparar diferentes técnicas de treinamento, incluindo estratégias de redução de ruído, \textit{Stochastic Weight Averaging} (SWA)~\cite{ref-izmailov2018} e \textit{focal loss}, analisando seus impactos no desempenho e na estabilidade dos modelos.
    \item Desenvolver e analisar modelos do tipo ensemble com o objetivo de aumentar a robustez e a estabilidade das predições.
    \item Realizar uma análise sistemática do desempenho das variantes da EfficientNet na tarefa de classificação dos graus ISUP.
    \item Avaliar o desempenho dos modelos propostos utilizando principalmente o \textit{Quadratic Weighted Kappa} (QWK)~\cite{ref-cohen1968}, métrica apropriada para problemas de classificação ordinal, em conjunto com a acurácia.
\end{itemize}

\section{Organização do Trabalho}
\label{sec:organizacao-do-trabalho}

Os demais capítulos deste trabalho foram organizados em:

\begin{itemize}
    \item O Capítulo~\ref{cap:trabalhos-relacionados} apresenta uma revisão dos trabalhos relacionados à classificação automatizada de graus ISUP em câncer de próstata.
    \item O Capítulo~\ref{cap:fundamentacao-teorica} trata dos conceitos fundamentais necessários para a construção desta pesquisa.
    \item O Capítulo~\ref{cap:materiais-metodo} descreve todas as etapas do método proposto, incluindo a implementação das arquiteturas EfficientNet, técnicas de treinamento avançadas e métodos de ensemble.
    \item Finalmente, o Capítulo~\ref{cap:conclusao} apresenta as considerações finais e sugestões de trabalhos futuros.
\end{itemize}